\chapter{Introdução}

Este capitulo deve apresentar o tema, a pergunta de pesquisa, a relevancia do estudo e a organizacao do trabalho.

\section{Contexto}

    A volatilidade dos preços de ativos financeiros é uma das variáveis mais relevantes em
finanças quantitativas. Ela representa a dispersão dos retornos de um ativo ao longo do tempo
e serve como insumo central para uma ampla quantidade de aplicações práticas, como o cálculo
do Value at Risk (VaR) e do Expected Shortfall (ES), medidas utilizadas por bancos e gestoras
de recursos para dimensionamento de reservas de capital, depende diretamente de estimativas
precisas da volatilidade. Da mesma forma, a precificação de opções e demais derivativos financeiros, 
fundamentada no modelo pioneiro de Black e Scholes (1973), e as estratégias de alocação de portfólio 
baseadas em controle de risco (volatility targeting) são diretamente sensíveis à qualidade dessas 
estimativas. No âmbito regulatório, o Acordo de Basileia III, estabelecido pelo Bank for International 
Settlements (BIS, 2019), vigente internacionalmente, exige que instituições financeiras utilizem 
o Expected Shortfall como métrica de risco de mercado, reforçando a importância de modelos de previsão de
volatilidade robustos.

    Entretanto, séries de retornos financeiros apresentam características que dificultam a modelagem
estatística convencional, incluindo heterocedasticidade condicional, clustering de volatilidade, distribuições com caudas pesadas
e evidências de assimetria na resposta da volatilidade a choques de mercado (Fama, 1965; Akgiray, 1989). Com isso, a família de modelos GARCH 
(Generalized Autoregressive Conditional Heteroskedasticity), introduzida por Engle (1982) e generalizada por Bollerslev 
(1986), foi desenvolvida para capturar essas características e permanece como referência consolidada na literatura acadêmica e 
na prática do setor financeiro. Além disso, extensões como o EGARCH (Nelson, 1991) incorporam o efeito alavancagem, em que choques 
negativos tendem a elevar a volatilidade de forma assimétrica em relação a choques positivos de mesma magnitude.

    Contudo, os modelos GARCH apresentam limitações relevantes. Por construção, os modelos 
são lineares na modelagem da variância condicional e têm dificuldade em capturar dinâmicas
não lineares complexas que frequentemente caracterizam mercados em
momentos de estresse. Por outro lado, o avanço dos algoritmos de machine learning nas últimas
décadas, em particular Random Forest, Gradient Boosting (XGBoost) e redes neurais
recorrentes do tipo Long Short-Term Memory (LSTM), abriram a possibilidade de abordagens
orientadas a dados com maior flexibilidade para identificar padrões não lineares. A questão
central que surge é se essa flexibilidade se traduz em ganho preditivo efetivo, especialmente
fora da amostra de treinamento.

    Dessa forma, o mercado acionário brasileiro constitui um ambiente particularmente relevante para a aplicação desses modelos. 
A Bolsa de Valores (B3) combina elevada volatilidade com forte sensibilidade a fatores externos, além da recorrência de choques 
políticos domésticos. Nesse contexto, torna-se relevante investigar se a maior flexibilidade dos algoritmos de machine learning, 
em comparação aos modelos econométricos tradicionais como o GARCH, se traduz em ganhos preditivos efetivos na previsão da volatilidade 
no mercado acionário brasileiro.

\section{Pergunta de pesquisa}

Modelos de machine learning (Random Forest, XGBoost e LSTM) superam modelos
GARCH (GARCH(1,1) e EGARCH) na previsão da volatilidade diária de ativos do mercado
acionário brasileiro?


\section{Objetivos}

\subsection{Objetivo Geral}

O objetivo do estudo é comparar o desempenho preditivo de modelos GARCH e de algoritmos de machine
learning na previsão da volatilidade diária de ativos do mercado acionário brasileiro, avaliando
em que condições cada abordagem apresenta maior acurácia.

\subsection{Objetivos Específicos}
Para alcançar o objetivo geral, definem-se os seguintes passos:
\begin{itemize}
    \item Estimar modelos GARCH(p,q) e EGARCH para a série de retornos do ativo
selecionado e obter previsões de volatilidade fora da amostra;
    \item Treinar e calibrar modelos de machine learning (Random Forest, XGBoost e LSTM)
para a previsão da volatilidade, com engenharia de features baseada em retornos passados e
indicadores de mercado;
    \item Analisar se a vantagem preditiva dos modelos se mantém em subperíodos de alta
volatilidade, como crises e choques macroeconômicos.
\end{itemize}

\section{Justificativa}

    A rigidez paramétrica dos modelos tradicionais de volatilidade constitui um desafio relevante na modelagem de séries financeiras, 
especialmente em contextos caracterizados por elevada instabilidade. No caso brasileiro, essa limitação torna-se ainda mais significativa,
uma vez que a dinâmica dos preços de ativos é fortemente influenciada por fatores externos, instabilidade macroeconômica e componentes 
comportamentais dos investidores.

    Nesse sentido, evidências empíricas sugerem que o comportamento dos agentes pode amplificar movimentos de mercado. Ferreira, Machado e Silva (2021) 
destacam que o pessimismo do investidor atua como um fator de propagação de incerteza, contribuindo para a formação de mispricings e para uma resposta 
assimétrica da volatilidade a novos choques de informação.

    Paralelamente, a literatura internacional aponta que modelos não lineares podem capturar dinâmicas de volatilidade não observadas por modelos econométricos 
tradicionais. Nesse contexto, Donaldson e Kamstra (1997) demonstram que redes neurais apresentam capacidade superior para modelar padrões complexos em séries financeiras,
sugerindo potencial ganho preditivo em relação à família GARCH.

    Apesar desses avanços, ainda não há consenso sobre a superioridade dessas abordagens, especialmente em mercados emergentes, onde características estruturais e 
institucionais podem afetar o desempenho dos modelos. Dessa forma, este trabalho contribui para a literatura ao comparar, no contexto do mercado acionário brasileiro, o 
desempenho preditivo de modelos GARCH e de algoritmos de machine learning na previsão da volatilidade, com foco na acurácia fora da amostra e em períodos de maior instabilidade.

\section{Estrutura do trabalho}

    Este trabalho está estruturado em seis capítulos, incluindo esta introdução. No segundo capítulo, apresenta-se a revisão de literatura,
que abrange a base teórica e as evidências empíricas pertinentes ao tema, fundamentando a evolução dos modelos de volatilidade e as aplicações 
de Machine Learning no setor financeiro. O terceiro capítulo detalha a metodologia, descrevendo o desenho  da pesquisa, os modelos propostos e 
as limitações do estudo. Em seguida, o quarto capítulo aborda a base de dados e a estratégia empírica, explicitando as fontes, a construção das 
variáveis e o tratamento de dados. Os resultados são discutidos no quinto capítulo, que compreende a análise descritiva, os resultados principais 
e testes de robustez, incluindo procedimentos de backtesting para validar a capacidade preditiva  e a estabilidade do modelo proposto. Por fim, o 
sexto capítulo apresenta a conclusão, sintetizando os achados, as contribuições do trabalho e as sugestões para pesquisas futuras. 
